% !PRINTED 27
\weberpage{27}
%<<CONT
\gesperrt{Der Theiler eines Productes zweier ganzer Functionen ist
gleich dem Product der Theiler beider Functionen.}

Denn sind $PA$ und $QB$ zwei ganze Functionen mit den Theilern $P$ und
$Q$, so sind $A$ und $B$ ursprüngliche Functionen. Also ist auch
$AB=C$ eine ursprüngliche Function und $PQ$ ist der Theiler der
Function $PA\cdot QB=PQC$.

Wir können dem Satze, insofern er sich auf Functionen einer
Veränderlichen bezieht, ohne seinen Inhalt wesentlich zu ändern,
folgende Fassung geben, in der er besonders nützlich ist. Sind
\[
 \varphi(x)=x^m+\alpha_1x^{m-1}+\alpha_2x^{m-2}+\cdots+\alpha_m,
\]
\[
 \psi(x)=x^n+\beta_1x^{n-1}+\beta_2x^{n-2}+\cdots+\beta_n
\]
\gesperrt{zwei ganze rationale Functionen, in denen die höchsten
Potenzen von $x$ den Coëfficienten $1$ haben, während die übrigen
Coëfficienten rationale Zahlen sind, so können in dem Product}
\[
 \varphi(x)\psi(x)=x^{m+n}+\gamma_1x^{m+n-1}
 +\gamma_2x^{m+n-2}+\cdots+\gamma_{m+n}
\]
\gesperrt{die Coëfficienten $\gamma$ nicht alle ganze Zahlen sein,
wenn die Coëfficienten $\alpha,\beta$ in $\varphi(x)$ und $\psi(x)$
nicht alle ganze Zahlen sind.}

Denn bezeichnen wir den kleinsten Hauptnenner der Coëfficienten
$\alpha$ von $\varphi$ mit $a_0$, der Coëfficienten $\beta$ von $\psi$
mit $b_0$, so sind $a_0\varphi(x)=A(x)$,
$b_0\psi(x)=B(x)$ primitive Functionen von $x$. Ihr Product
$a_0b_0\varphi(x)\psi(x)$ hätte, wenn die Coëfficienten $\gamma$ ganze
Zahlen wären, den Theiler $a_0b_0$ und wäre also, wenn $a_0$ und $b_0$
nicht beide gleich $1$ wären, nicht primitiv; dies aber wäre ein
Widerspruch mit dem oben bewiesenen Satze.

\begin{center}
{\large \S.\ 3.}

\vspace{0.5em}
{\large\gesperrt{Division.}}
\end{center}

Es seien, wie bisher,
\begin{equation}
\begin{aligned}
 A&=A(x)=a_0x^m+a_1x^{m-1}+\cdots,\\
 B&=B(x)=b_0x^n+b_1x^{n-1}+\cdots
\end{aligned}
\tag{1}
\end{equation}
zwei ganze rationale Functionen von $x$; es soll aber jetzt voraus-
%CONT>>
